Source of systematic uncertainty of the methods presented above
potentially arises from the different composition of background processes 
($DY$, $t \bar{t}$, $WZ$, $Z\gamma$) in the region where we measure and 
where we apply the fake ratios.
OS method corrects for the resulting bias via the ``3P1F component"
of its prediction. SS method corrects explicitly the electron fake rates
for the correct fraction of photon conversions, but no attempt is made 
to correct the muon fake rates.
The closure tests presented here are used to assess a possible residual
bias in the methods.

\paragraph{Statistics in 4l Control Sample}
The limited size of the samples in the control regions where we measure and where we apply the fake ratios is the source of the statistical uncertainties of the method. The dominating statistical uncertainty is driven by the number of events in the control region and is typically in the range of 1-10\%.

\paragraph{Sensitivity of Fake Ratios to Background Composition}
Compositions of reducible background processes ($DY$, $t \bar{t}$, $WZ$, $Z\gamma^{(*)}$) in the region where we measure and where we apply the fake ratio method are typically not the same. This is the main source of the systematic uncertainty of the fake ratio method.

This uncertainty can be estimated by measuring the fake ratios for individual background processes in the $Z+1L$ region in simulation. The weighted average of these individual fake ratios is the fake ratio that we measure in this sample (in simulation). The exact composition of the background processes in the 2P+2F region where we plan to apply the fake ratios can be determined from simulation, and one can reweigh the individual fake ratios according to the 2P+2F composition. The difference between the reweighed fake ratio and the average one can be used as a measure of the uncertainty on the measurement of the fake ratios. The fake ratios for individual processes, the average fake ratio and the reweighed fake ratio determined by simulation are shown in table below. The effect of this systematic uncertainty is propagated to the final estimates, and it amounts to about 32\% for $4e$, 33\% for $2e2\mu$ and 35\% for $4\mu$ final state. 

% \begin{table}[h]
% \scriptsize
%     \centering
%     \caption{
%     The fake ratios for individual background processes, the average fake ratio and the fake ratio reweighed according to the composition of backgrounds in 2P+2F region.}  
    
%     \begin{tabular}[!htb]{| l | c | c | c | c | c | c |} \hline
% Sample				& Light Jets		  & HF Jets& From $\gamma$ conv	    & average (in $Z+1\ell$) & reweighed (2P2F)	& uncertainty \\ \hline \hline
% electron fake ratio	& $0.012^{+0.001}_{-0.001}$ 	& $0.115^{+0.005}_{-0.006}$  & $0.293^{+0.043}_{-0.043}$  &$0.021^{+0.001}_{-0.001}$& $0.024^{+0.004}_{-0.004}$& 15\%    \\ 
% muon fake ratio 	& $0.057^{+0.003}_{-0.003}$ 	& $0.225^{+0.008}_{-0.008}$ & $0.003^{+0.455}_{-0.003}$ &$0.120^{+0.004}_{-0.004}$& $0.105^{+0.020}_{-0.017}$ & 13\% \\  \hline
% 	\end{tabular}
%     \label{tab:uncertFR}
% \end{table}

\paragraph{Shape Uncertainty}
In order to estimate the uncertainty on the $m_{4l}$ shape we have looked at the differences between the shapes of predicted background distributions for all three channels, and between both of the two methods. The envelope of differences between these shapes of distributions is used as an estimate of the shape uncertainty. The uncertainty is estimated to be roughly in the range of 5\% - 15\%. Since the difference of the shapes slowly varies with $m_{4l}$, it is taken as a constant versus $m_{4l}$  and is absorbed in the much larger uncertainty on the predicted yield of backgrounds. 
%The shapes of predicted background $m_{4l}$ distributions for the three channels are shown in Figure~\ref{fig:SR_CombinedShapes} (left).	

